% Options for packages loaded elsewhere
\PassOptionsToPackage{unicode}{hyperref}
\PassOptionsToPackage{hyphens}{url}
\documentclass[
]{article}
\usepackage{xcolor}
\usepackage{amsmath,amssymb}
\setcounter{secnumdepth}{-\maxdimen} % remove section numbering
\usepackage{iftex}
\ifPDFTeX
  \usepackage[T1]{fontenc}
  \usepackage[utf8]{inputenc}
  \usepackage{textcomp} % provide euro and other symbols
\else % if luatex or xetex
  \usepackage{unicode-math} % this also loads fontspec
  \defaultfontfeatures{Scale=MatchLowercase}
  \defaultfontfeatures[\rmfamily]{Ligatures=TeX,Scale=1}
\fi
\usepackage{lmodern}
\ifPDFTeX\else
  % xetex/luatex font selection
\fi
% Use upquote if available, for straight quotes in verbatim environments
\IfFileExists{upquote.sty}{\usepackage{upquote}}{}
\IfFileExists{microtype.sty}{% use microtype if available
  \usepackage[]{microtype}
  \UseMicrotypeSet[protrusion]{basicmath} % disable protrusion for tt fonts
}{}
\makeatletter
\@ifundefined{KOMAClassName}{% if non-KOMA class
  \IfFileExists{parskip.sty}{%
    \usepackage{parskip}
  }{% else
    \setlength{\parindent}{0pt}
    \setlength{\parskip}{6pt plus 2pt minus 1pt}}
}{% if KOMA class
  \KOMAoptions{parskip=half}}
\makeatother
\setlength{\emergencystretch}{3em} % prevent overfull lines
\providecommand{\tightlist}{%
  \setlength{\itemsep}{0pt}\setlength{\parskip}{0pt}}
\usepackage{bookmark}
\IfFileExists{xurl.sty}{\usepackage{xurl}}{} % add URL line breaks if available
\urlstyle{same}
% fallback for those not using the hyperref driver hyperxmp:
\makeatletter
\@ifundefined{xmpquote}{\newcommand{\xmpquote}[1]{#1}}{}
\makeatother
\hypersetup{
  hidelinks,
  pdfcreator={LaTeX via pandoc}}

\author{}
\date{}

\begin{document}

\section{Prediction-guided mixed-model GWAS refines maize oil-trait
candidate intervals in
ZEAMAP}\label{prediction-guided-mixed-model-gwas-refines-maize-oil-trait-candidate-intervals-in-zeamap}

\subsection{Title Page}\label{title-page}

Title: Prediction-guided mixed-model GWAS refines maize oil-trait
candidate intervals in ZEAMAP

Running title: ZEAMAP oil-trait prediction and GWAS

Authors: ×××

Affiliations: ×××

Corresponding author: ×××

ORCID/email: ×××

Author confirmation: ×××

\subsection{Core Ideas}\label{core-ideas}

\begin{itemize}
\tightlist
\item
  ZEAMAP processed data were harmonized into a conservative
  accession-level analysis set.
\item
  Oil and fatty-acid traits were more predictable than other trait
  families under regularized genotype-based models.
\item
  Mixed-model GWAS was essential: covariate-only testing was inflated,
  whereas GEMMA controlled genomic inflation.
\item
  Recurrent chr6 and chr9 intervals provide the clearest fatty-acid
  candidate signals.
\item
  The study prioritizes candidate intervals for follow-up and does not
  claim causal variants or experimentally validated genes.
\end{itemize}

\subsection{Abstract}\label{abstract}

Maize kernel oil content and fatty-acid composition are important
seed-quality traits, but association analyses in diverse maize panels
must distinguish true local signals from population structure and
relatedness. We harmonized public ZEAMAP processed data into a
conservative accession-level analysis set containing 461 accessions,
199,856 filtered SNPs and 318 numeric phenotype or metabolite traits.
Repeated genotype-to-phenotype benchmarking identified 66 robustly
predictable traits; oil traits formed the strongest family, with a
genotype plus population ridge model reaching median Pearson/R2 of
0.597/0.338 across robust oil traits. This prediction result motivated
focused GWAS for 10 high-priority oil traits. A covariate-only scan was
strongly inflated, whereas GEMMA mixed linear models with
genotype-derived kinship controlled lambda GC to 0.984-1.018 (median
0.998). We annotated and ranked 184 candidate loci and 147 candidate
genes after excluding nominal-only signals. The strongest evidence
converged on a recurrent chromosome 6 linoleic acid1-region interval and
a chromosome 9 C16:0-associated interval containing nearby acyl-ACP
thioesterase annotation. These results provide a reproducible framework
for moving from public maize processed data to calibrated oil-trait
candidate intervals while keeping causal claims appropriately bounded.

\subsection{Keywords}\label{keywords}

maize; ZEAMAP; kernel oil; fatty-acid composition; genomic prediction;
mixed-model GWAS; GEMMA; candidate interval

\subsection{Introduction}\label{introduction}

Kernel oil concentration and fatty-acid composition influence maize
grain quality, nutritional value and industrial use. The traits are also
biologically informative because they connect quantitative genetic
variation with fatty-acid metabolism, seed development and carbon
allocation. Previous maize studies have identified oil-content and
fatty-acid loci, including signals related to linoleic-acid composition
and DGAT-mediated oil accumulation (Alrefai et al., 1995; Zheng et al.,
2008; Li et al., 2013; Zhang et al., 2023). Yet public diversity-panel
data still require careful handling before these biological signals can
be interpreted.

ZEAMAP provides a broad public resource for maize multi-omics and trait
data (Gui et al., 2020). The practical challenge is that processed
tables do not automatically form a single analysis matrix. Accessions,
modalities and biological units must be reconciled before modelling. In
particular, reference or tissue-level expression resources cannot be
treated as accession-level expression measurements, and partially paired
methylation data should not be used as if they covered the full panel.
We therefore built a restrained accession-level analysis set in which
genotype, population information and phenotype/metabolite traits were
aligned before modelling.

The second challenge is statistical. A panel of 461 accessions and
nearly 200,000 SNPs is large enough for genome-wide association and
regularized prediction, but not large enough to justify unconstrained
high-capacity prediction models. A model can appear sophisticated while
learning population structure or noise. We therefore used regularized
models and repeated evaluation to identify trait families with stable
predictive signal, then used that result to prioritize a focused GWAS
layer.

The third challenge is association calibration. Diversity-panel GWAS is
vulnerable to residual structure. A candidate-gene story built on
inflated tests is weak even when the gene annotation is appealing. For
this reason we explicitly treated covariate-only GWAS as a diagnostic
layer and used GEMMA mixed linear models as the manuscript-level
association test (Zhou \& Stephens, 2012; Zhou \& Stephens, 2014). Lead
loci were then interpreted through a conservative evidence hierarchy
combining statistical support, recurrence across oil traits,
prediction-overlap information and local gene annotation from B73
reference resources (Jiao et al., 2017; Yates et al., 2022).

\subsection{Results}\label{results}

\subsubsection{Accession harmonization produced a restrained ZEAMAP
analysis
set}\label{accession-harmonization-produced-a-restrained-zeamap-analysis-set}

After harmonizing processed ZEAMAP genotype, population and
phenotype/metabolite tables, the analysis set retained 461 accessions
with the required paired information, 199,856 filtered biallelic SNPs
and 318 numeric traits. Methylation coverage was available for 236
accessions and was tracked as a partially paired modality rather than
merged into the primary model. Expression resources available in the
working data were not accession-matched to the diversity panel and were
therefore excluded from genotype-to-phenotype prediction.

This filtering reduced the apparent size of the resource but increased
interpretability. Every primary prediction and GWAS result was based on
the same accession-level unit. This matters because oil-trait signals in
a structured maize panel can be confounded if sample identity or
modality pairing is loose.

\subsubsection{Oil traits showed the strongest and most stable
prediction
signal}\label{oil-traits-showed-the-strongest-and-most-stable-prediction-signal}

Repeated prediction benchmarking showed that oil traits were the
clearest target family. Across selected trait families, oil traits had
the highest median Pearson correlation (0.597) and median R2 (0.338),
with the best oil trait reaching Pearson/R2 of 0.924/0.827. Among the
robust final traits, 29 were oil-related. The best-performing traits
included total oil content and major fatty-acid components, indicating
that the genotype and population features captured biologically
meaningful variation rather than only broad agronomic structure.

Model comparison also supported a conservative modelling choice. Ridge
and ElasticNet were more stable than a small multilayer perceptron under
the available sample size. The result is not a claim against neural
networks in maize genetics; it shows that, for this data scale and
pairing structure, regularized models provide a stronger basis for trait
prioritization and biological follow-up.

\subsubsection{Methylation was informative only as an auxiliary
layer}\label{methylation-was-informative-only-as-an-auxiliary-layer}

Methylation was evaluated in three forms: global accession summaries,
gene/promoter/cis-window principal components and sparse gene-window
features. The global features showed little broad improvement,
gene-level components provided limited auxiliary signal and sparse
methylation features were unstable in the smaller 236-accession subset.
We therefore did not include methylation in the primary prediction
model. This decision prevents a weak multi-omics claim and keeps the
main analysis anchored in the fully paired genotype and trait data.

\subsubsection{GEMMA controlled inflation that persisted in
covariate-only
GWAS}\label{gemma-controlled-inflation-that-persisted-in-covariate-only-gwas}

Prediction results directed the association analysis toward 10
high-priority oil traits. The covariate-only scan was useful as a
warning: genomic inflation remained high despite population covariates.
In contrast, GEMMA mixed linear models using genotype-derived kinship
controlled lambda GC to 0.984-1.018, with a median of 0.998. This
calibration is central to the study because it supports interpretation
of local peaks as candidate intervals rather than artifacts of residual
relatedness.

Across the 10 oil traits, GEMMA tested 440 non-missing accessions and
199,856 SNPs per trait. Bonferroni-significant hits ranged from 1 to 21
per trait, with 126 Bonferroni-level hits and 338 FDR-level hits in
total. The strongest signals occurred for very-long-chain fatty-acid
composition, C18 fatty-acid traits and C16:0. Trait-level Manhattan and
QQ plots were retained as diagnostic figures, while the main text
focuses on calibrated summary statistics and prioritized intervals.

\subsubsection{Candidate-locus ranking separated direct lipid candidates
from broader regulatory
regions}\label{candidate-locus-ranking-separated-direct-lipid-candidates-from-broader-regulatory-regions}

GEMMA lead signals were merged into physical candidate intervals and
annotated with B73 RefGen\_v4 gene models. After removing nominal-only
signals, the candidate set contained 184 loci and 147 candidate genes.
Of these loci, 63 also overlapped ridge-attribution evidence and 11
carried fatty-acid or lipid-related annotation. This layered ranking
separated two types of evidence: direct lipid-metabolism candidates and
recurrent but functionally broader regions.

The strongest interval was R01\_Zm00001d036982 on chromosome 6. It was
recurrent across seven oil traits, reached best P = 2.35e-25 and
contained the B73 RefGen\_v4 gene Zm00001d036982, which maps to the
current B73 v5 gene Zm00001eb277490. Local annotation and external xrefs
support lipid acyltransferase/DGAT-like or linoleic-acid biology. This
region therefore has the strongest combination of recurrence,
statistical support, prediction overlap and prior fatty-acid evidence.

The chr9 C16:0 interval provides a different but biologically important
signal. The lead interval around Zm00001d045383 reached best P =
7.76e-17 for C16:0. The lead gene maps to a DXS/isoprenoid-related
annotation, while the nearby Zm00001d045387/Zm00001eb377350 gene carries
acyl-ACP hydrolase, palmitoyl-ACP thioesterase and fatty-acid
biosynthesis evidence. This makes the region a high-priority saturated
fatty-acid candidate interval, but not yet a fine-mapped FatB causal
gene.

Other tier-1 intervals on chromosomes 1, 4 and 8 were recurrent and
ridge-supported but had broader annotations, including MYB-domain,
protein-transport and tetratricopeptide-repeat genes. These regions are
valuable because recurrence across oil traits suggests reproducible
association structure. However, their current annotations support
candidate-interval reporting rather than direct oil-biosynthetic claims.

\subsubsection{Expanded figure and table package clarifies the evidence
hierarchy}\label{expanded-figure-and-table-package-clarifies-the-evidence-hierarchy}

The revised figure set contains four main figures. Figure 1 summarizes
data harmonization and prediction benchmarking. Figure 2 presents the
mixed-model GWAS calibration and candidate-locus summary. Figure 3 shows
the chr6 and chr9 regional intervals. Figure 4 integrates trait-family
predictability, GEMMA calibration, top-region association strength and
evidence classes across prioritized loci. The table set contains one
main evidence table for the eight top regional loci and four
supplementary tables covering the full candidate-locus set, tier-1 loci,
external annotation hardening and column definitions. This structure
keeps the main text readable while giving reviewers enough information
to inspect the full locus-ranking logic.

\subsection{Discussion}\label{discussion}

This study provides a calibrated route from public ZEAMAP processed data
to maize oil-trait candidate intervals. The most important result is not
simply that several oil-trait associations were detected. Rather, the
analysis shows why these associations are interpretable: the input
accessions were harmonized, oil traits were prioritized by repeated
prediction, covariate-only GWAS was shown to be inflated and GEMMA mixed
models brought genomic inflation close to one.

The prediction results help explain why the GWAS focus on oil traits is
justified. Oil and fatty-acid traits were not chosen only because they
are biologically attractive. They were the most consistently predictable
trait family in the analysis set, with strong performance for total oil
and major fatty-acid components. This suggests that the available SNP
features capture a meaningful part of the genetic architecture for these
traits. In practical breeding terms, the result supports oil traits as a
near-term target for regularized genomic prediction in this
ZEAMAP-derived panel, while more weakly predictable metabolite or
amino-acid traits may require larger sample sizes, better environmental
metadata or different feature representations.

The contrast between covariate-only and mixed-model GWAS is equally
important. In a structured maize panel, population covariates alone were
not sufficient. Without the kinship term, the association scan produced
inflation severe enough to undermine candidate interpretation. GEMMA did
not merely improve a technical diagnostic; it changed the credibility of
the biological conclusions. The near-one lambda values across all 10 oil
traits indicate that the final candidate intervals are less likely to be
dominated by residual relatedness.

The chr6 interval is the main biological anchor of the manuscript. Its
strength comes from convergence: multiple oil traits point to the same
region, the association is highly significant, ridge-derived evidence
supports the region and local annotation connects it to linoleic-acid or
lipid-acyltransferase biology. Previous maize oil and fatty-acid studies
provide a plausible context for this signal (Alrefai et al., 1995; Zheng
et al., 2008; Li et al., 2013; Zhang et al., 2023). The appropriate
claim is therefore strong but bounded: this is the leading recurrent
fatty-acid candidate interval in the present analysis, not a validated
causal gene.

The chr9 C16:0 interval is more trait-specific and illustrates why
gene-level interpretation must be careful. The lead gene itself is not
annotated as FatB. The stronger fatty-acid interpretation comes from a
nearby acyl-ACP thioesterase candidate, a pathway class directly related
to saturated fatty-acid composition. This supports follow-up of the
interval, especially for C16:0, but the result should not be overstated
as a confirmed causal FatB allele.

The broader set of recurrent loci may be useful for breeding and
functional prioritization. Some intervals contain regulatory or
transport-related candidates rather than canonical oil-pathway genes.
Such genes could still affect seed composition through developmental
timing, transport, organelle function or resource allocation. However,
these hypotheses need external expression evidence, fine mapping or
functional assays. The revised evidence hierarchy makes this distinction
explicit: lipid-annotated intervals are discussed as stronger
candidates, while broader recurrent regions are retained as
hypothesis-generating intervals.

Several limitations define the next experiments. First, the analysis set
remains modest for high-dimensional prediction, so complex models should
wait for larger paired datasets. Second, methylation was only partially
paired and cannot yet support a strong multi-omics mechanism. Third, the
GWAS resolution is limited by LD, marker density and accession sample
size. Fourth, candidate genes were not experimentally validated.
Stronger follow-up would include independent panels, seed-specific
expression evidence, haplotype analysis around chr6 and chr9, and
functional tests for the leading lipid-related genes.

Overall, the study is best positioned as a plant genomics resource and
candidate-prioritization paper. It contributes a reproducible
accession-level analysis, a statistically calibrated oil-trait GWAS
layer and a ranked set of candidate intervals for maize oil and
fatty-acid composition. The claims are intentionally conservative, which
should make the manuscript more defensible in review.

\subsection{Materials and Methods}\label{materials-and-methods}

\subsubsection{ZEAMAP data and accession
matching}\label{zeamap-data-and-accession-matching}

Processed ZEAMAP public data were obtained from CNGBdb project
CNP0001565 (Gui et al., 2020). Genotype, population, phenotype and
metabolome tables were matched at accession level. Accessions were
retained when genotype, population information and at least one
phenotype or metabolome measurement were present. SNPs were filtered to
common biallelic markers with sufficient sample coverage, yielding
199,856 variants for modelling and association analysis.

\subsubsection{Prediction analysis}\label{prediction-analysis}

Genotype-to-phenotype prediction was evaluated using repeated
train/validation/test splits. Candidate models included genotype plus
population ridge regression, genotype plus population ElasticNet,
population-only ridge regression and a small multilayer perceptron.
Traits were retained for final reporting only when performance was
positive and stable across seeds. Pearson correlation and R2 on held-out
samples were used as the main performance metrics.

\subsubsection{Methylation assessment}\label{methylation-assessment}

Methylation information was evaluated as an auxiliary modality in the
subset of accessions with coverage. We tested global methylation
summaries, gene/promoter/cis-window components and sparse gene-window
features. Because the methylation subset was substantially smaller than
the full analysis set and sparse features were unstable, methylation was
not used in the primary prediction model.

\subsubsection{Mixed-model GWAS}\label{mixed-model-gwas}

Ten high-priority oil traits were selected from the robust prediction
results. Covariate-only GWAS was used as a diagnostic comparison. The
primary association analysis used GEMMA mixed linear models with
genotype-derived kinship and population covariates (Zhou \& Stephens,
2012; Zhou \& Stephens, 2014). Likelihood-ratio p-values were used for
locus ranking. The Bonferroni threshold was 0.05 divided by 199,856
tested variants; FDR summaries were used for candidate triage rather
than causal claims (Benjamini \& Hochberg, 1995).

\subsubsection{Candidate-locus annotation and
ranking}\label{candidate-locus-annotation-and-ranking}

Lead SNPs were merged into physical intervals and annotated with B73
RefGen\_v4 gene models (Jiao et al., 2017). Candidate intervals were
ranked by association strength, recurrence across oil traits, overlap
with prediction-derived evidence, local gene annotation and external
gene-model support from plant genome resources (Yates et al., 2022).
Intervals were classified as direct lipid candidates, indirect
regulatory/transport candidates or unresolved recurrent candidates.
Candidate genes were treated as hypotheses unless supported by
independent functional evidence.

\subsubsection{Software and
reproducibility}\label{software-and-reproducibility}

Analyses used Python scientific-computing libraries and GEMMA for
mixed-model association (Pedregosa et al., 2011; Harris et al., 2020;
Virtanen et al., 2020; McKinney, 2010; Hunter, 2007; Zhou \& Stephens,
2012). The repository includes the scripts and summary tables needed to
reproduce the accession harmonization, prediction benchmark, GWAS
summaries, candidate-locus ranking and manuscript figures.

\subsection{Data Availability
Statement}\label{data-availability-statement}

The study reanalyzes public ZEAMAP processed data from CNGBdb project
CNP0001565. The project repository contains analysis scripts, summary
tables, candidate-locus annotations and generated figures. Large public
raw inputs and large intermediate matrices are not included in the
repository; their source and regeneration paths are documented with the
analysis files.

\subsection{Code Availability
Statement}\label{code-availability-statement}

Code used for dataset construction, prediction benchmarking, methylation
assessment, GWAS summarization, candidate-locus annotation and figure
preparation is available in the project repository under
\texttt{scripts/}.

\subsection{Author Contributions}\label{author-contributions}

×××

\subsection{Funding}\label{funding}

×××

\subsection{Acknowledgements}\label{acknowledgements}

×××

\subsection{Competing Interests}\label{competing-interests}

×××

\subsection{Ethics Statement}\label{ethics-statement}

No human participants or animal subjects were used. The study reanalyzes
public plant genomics data.

\subsection{Figure Legends}\label{figure-legends}

Figure 1. ZEAMAP accession harmonization and prediction benchmark. The
figure summarizes accession matching, retained data scale, model
comparison and trait-family prediction performance.

Figure 2. GEMMA mixed-model GWAS calibration and candidate-locus
summary. The figure shows that GEMMA controlled genomic inflation and
summarizes the filtered candidate-locus set.

Figure 3. Regional evidence for chr6 and chr9 oil-trait candidate
intervals. Panels highlight the recurrent chr6 linoleic acid1-region
interval and the chr9 C16:0-associated acyl-ACP thioesterase candidate
interval.

Figure 4. Integrated prediction and GWAS evidence. Panels summarize
trait-family predictability, GEMMA calibration across oil traits,
association strength of top regional loci and evidence classes across
prioritized regions.

\subsection{Table Legends}\label{table-legends}

Table 1. Prioritized regional candidate loci for maize oil traits. Eight
top loci are summarized by recurrence, association strength, prediction
overlap, functional annotation and claim boundary.

Supplementary Table 1. Manuscript candidate loci from oil-trait GEMMA
mixed-model GWAS.

Supplementary Table 2. Tier-1 candidate loci used for main-text
interpretation.

Supplementary Table 3. External annotation review for top regional loci.

Supplementary Table 4. Column dictionary for supplementary tables.

\subsection{References}\label{references}

Alrefai, R., Berke, T. G., \& Rocheford, T. R. (1995). Quantitative
trait locus analysis of fatty acid concentrations in maize.
\emph{Genome}, \emph{38}, 894-901. https://doi.org/10.1139/g95-118

Benjamini, Y., \& Hochberg, Y. (1995). Controlling the false discovery
rate: A practical and powerful approach to multiple testing.
\emph{Journal of the Royal Statistical Society: Series B}, \emph{57},
289-300. https://doi.org/10.1111/j.2517-6161.1995.tb02031.x

Bonaventure, G., Salas, J. J., Pollard, M. R., \& Ohlrogge, J. B.
(2003). Disruption of the FATB gene in Arabidopsis demonstrates an
essential role of saturated fatty acids in plant growth. \emph{The Plant
Cell}, \emph{15}, 1020-1033. https://doi.org/10.1105/tpc.008946

Cook, J. P., McMullen, M. D., Holland, J. B., Tian, F., Bradbury, P.,
Ross-Ibarra, J., Buckler, E. S., \& Flint-Garcia, S. A. (2012). Genetic
architecture of maize kernel composition in the nested association
mapping and inbred association panels. \emph{Plant Physiology},
\emph{158}, 824-834. https://doi.org/10.1104/pp.111.185033

Gui, S., Yang, L., Li, J., Luo, J., Xu, X., Yuan, J., Chen, L., Li, W.,
Yang, X., \& Wang, J. (2020). ZEAMAP, a comprehensive database adapted
to the maize multi-omics era. \emph{iScience}, \emph{23}, 101241.
https://doi.org/10.1016/j.isci.2020.101241

Harris, C. R., Millman, K. J., van der Walt, S. J., Gommers, R.,
Virtanen, P., Cournapeau, D., Wieser, E., Taylor, J., Berg, S., Smith,
N. J., Kern, R., Picus, M., Hoyer, S., van Kerkwijk, M. H., Brett, M.,
Haldane, A., Del Rio, J. F., Wiebe, M., Peterson, P., \ldots{} Oliphant,
T. E. (2020). Array programming with NumPy. \emph{Nature}, \emph{585},
357-362. https://doi.org/10.1038/s41586-020-2649-2

Hunter, J. D. (2007). Matplotlib: A 2D graphics environment.
\emph{Computing in Science \& Engineering}, \emph{9}, 90-95.
https://doi.org/10.1109/MCSE.2007.55

Jiao, Y., Peluso, P., Shi, J., Liang, T., Stitzer, M. C., Wang, B.,
Campbell, M. S., Stein, J. C., Wei, X., Chin, C. S., Guill, K.,
Regulski, M., Kumari, S., Olson, A., Gent, J., Schneider, K. L.,
Wolfgruber, T. K., May, M. R., Springer, N. M., \ldots{} Ware, D.
(2017). Improved maize reference genome with single-molecule
technologies. \emph{Nature}, \emph{546}, 524-527.
https://doi.org/10.1038/nature22971

Katral, A., Ahirwar, R. N., Gazala, P., Jaiswal, S. K., Sachan, M.,
Barupal, T., Hossain, F., Muthusamy, V., Saripalli, G., Mishra, S. J.,
Gupta, H. S., \& Chhabra, R. (2022). Allelic variation in Zmfatb gene
defines variability for fatty acids composition among diverse maize
genotypes. \emph{Frontiers in Nutrition}, \emph{9}, 845255.
https://doi.org/10.3389/fnut.2022.845255

Li, H., Peng, Z., Yang, X., Wang, W., Fu, J., Wang, J., Han, Y., Chai,
Y., Guo, T., Yang, N., Liu, J., Warburton, M. L., Cheng, Y., Hao, X.,
Zhang, P., Zhao, J., Liu, Y., Wang, G., Li, J., \& Yan, J. (2013).
Genome-wide association study dissects the genetic architecture of oil
biosynthesis in maize kernels. \emph{Nature Genetics}, \emph{45}, 43-50.
https://doi.org/10.1038/ng.2484

McKinney, W. (2010). Data structures for statistical computing in
Python. In \emph{Proceedings of the 9th Python in Science Conference}
(pp.~56-61). https://doi.org/10.25080/Majora-92bf1922-00a

Pedregosa, F., Varoquaux, G., Gramfort, A., Michel, V., Thirion, B.,
Grisel, O., Blondel, M., Prettenhofer, P., Weiss, R., Dubourg, V.,
Vanderplas, J., Passos, A., Cournapeau, D., Brucher, M., Perrot, M., \&
Duchesnay, E. (2011). Scikit-learn: Machine learning in Python.
\emph{Journal of Machine Learning Research}, \emph{12}, 2825-2830.
https://jmlr.org/papers/v12/pedregosa11a.html

Virtanen, P., Gommers, R., Oliphant, T. E., Haberland, M., Reddy, T.,
Cournapeau, D., Burovski, E., Peterson, P., Weckesser, W., Bright, J.,
van der Walt, S. J., Brett, M., Wilson, J., Millman, K. J., Mayorov, N.,
Nelson, A. R. J., Jones, E., Kern, R., Larson, E., \ldots{} van
Mulbregt, P. (2020). SciPy 1.0: Fundamental algorithms for scientific
computing in Python. \emph{Nature Methods}, \emph{17}, 261-272.
https://doi.org/10.1038/s41592-019-0686-2

Yates, A. D., Allen, J., Amode, R. M., Azov, A. G., Barba, M., Becerra,
A., Bhai, J., Campbell, L. I., Carbajo Martinez, M., Chakiachvili, M.,
Chougule, K., Christensen, M., Contreras-Moreira, B., Cuzick, A., Da Rin
Fioretto, L., Davis, P., De Silva, N., Diamantakis, S., Dyer, S. C.,
\ldots{} Flicek, P. (2022). Ensembl Genomes 2022: An expanding genome
resource for non-vertebrates. \emph{Nucleic Acids Research}, \emph{50},
D996-D1003. https://doi.org/10.1093/nar/gkab1007

Zhang, C., Wang, J., Wang, L., Zhang, W., Liu, J., Li, H., Zhang, Y.,
Zhang, J., Yang, X., \& Yan, J. (2023). Genetic dissection of QTLs for
oil content in four maize DH populations. \emph{Frontiers in Plant
Science}, \emph{14}, 1174985. https://doi.org/10.3389/fpls.2023.1174985

Zheng, P., Allen, W. B., Roesler, K., Williams, M. E., Zhang, S., Li,
J., Glassman, K., Ranch, J., Nubel, D., Solawetz, W., Bhattramakki, D.,
Llaca, V., Deschamps, S., Zhong, G. Y., Tarczynski, M. C., \& Shen, B.
(2008). A phenylalanine in DGAT is a key determinant of oil content and
composition in maize. \emph{Nature Genetics}, \emph{40}, 367-372.
https://doi.org/10.1038/ng.85

Zhou, X., \& Stephens, M. (2012). Genome-wide efficient mixed-model
analysis for association studies. \emph{Nature Genetics}, \emph{44},
821-824. https://doi.org/10.1038/ng.2310

Zhou, X., \& Stephens, M. (2014). Efficient multivariate linear mixed
model algorithms for genome-wide association studies. \emph{Nature
Methods}, \emph{11}, 407-409. https://doi.org/10.1038/nmeth.2848

\end{document}
